% Editable source of truth. Compile from docs/ with: make report
% Update the evidence date and bibliography when changing scientific claims.
\documentclass[UTF8,fontset=fandol,zihao=5,a4paper]{ctexart}
\usepackage[margin=22mm,top=23mm,bottom=23mm,headheight=15pt]{geometry}
\usepackage{amsmath,amssymb,booktabs,tabularx,array,enumitem}
\usepackage[table]{xcolor}
\usepackage{graphicx,tikz,caption,fancyhdr,hyperref,xurl}
\usetikzlibrary{arrows.meta,positioning,calc,fit,backgrounds}
\definecolor{ink}{HTML}{18354A}
\definecolor{teal}{HTML}{087E8B}
\definecolor{gold}{HTML}{B67827}
\definecolor{pale}{HTML}{EDF5F7}
\definecolor{muted}{HTML}{586975}
\hypersetup{colorlinks=true,linkcolor=teal,citecolor=teal,urlcolor=teal,
 pdftitle={玄女：全国年度5米地理嵌入技术报告},pdfauthor={heyuhang},bookmarksnumbered=true}
\pagestyle{fancy}\fancyhf{}
\fancyhead[L]{\small\color{muted}玄女 / 年度 5m 地理嵌入}
\fancyhead[R]{\small\color{muted}技术方案与实施依据 · v3.1}
\fancyfoot[C]{\small\thepage}
\renewcommand{\headrulewidth}{0.3pt}
\ctexset{section={format=\Large\bfseries\color{ink}},subsection={format=\large\bfseries\color{ink}}}
\setlength{\parskip}{4pt}
\setlength{\emergencystretch}{3em}
\setlist{nosep,leftmargin=1.6em,topsep=4pt}
\renewcommand{\arraystretch}{1.25}
\captionsetup{font=small,labelfont=bf,justification=raggedright,singlelinecheck=false}
\newcolumntype{Y}{>{\raggedright\arraybackslash}X}
\newcommand{\status}[1]{\textcolor{teal}{\textbf{#1}}}
\newcommand{\note}[1]{\par\smallskip\noindent\colorbox{pale}{\parbox{0.95\linewidth}{\small #1}}\par\smallskip}
\newcommand{\source}[1]{\par{\footnotesize\color{muted}依据：#1}\par}
\newcommand{\shape}[1]{\ensuremath{#1}}
\newcommand{\reportpage}{\clearpage}
\tikzset{box/.style={draw=teal!65,fill=pale,rounded corners=2pt,align=center,
 font=\small,inner sep=6pt},goldbox/.style={box,draw=gold,fill=gold!7},
 arr/.style={-{Latex[length=2mm]},draw=ink,line width=.65pt},
 labeltext/.style={font=\footnotesize,text=muted,align=center}}

\begin{document}
\hypersetup{pageanchor=false}
\begin{titlepage}
\thispagestyle{empty}
\setlength{\parindent}{0pt}
\setlength{\parskip}{2pt}
\begin{tikzpicture}[remember picture,overlay]
\fill[ink] (current page.north west) rectangle ([yshift=-4mm]current page.north east);
\begin{scope}[shift={([xshift=-56mm,yshift=-24mm]current page.north east)}]
\draw[step=6mm,teal!13,line width=.3pt] (0,0) grid (30mm,-30mm);
\fill[teal!10] (6mm,-6mm) rectangle (12mm,-12mm);
\fill[teal!20] (18mm,-18mm) rectangle (24mm,-24mm);
\draw[teal!35,line width=.7pt] (12mm,-12mm) rectangle (24mm,-24mm);
\end{scope}
\end{tikzpicture}
\vspace*{5mm}
{\sffamily\bfseries\large\color{teal}XUANNV EMBEDDING\par}
\vspace{39mm}
{\fontsize{30}{40}\selectfont\bfseries\color{ink}全国年度 5 米\par 地理嵌入技术报告\par}
\vspace{7mm}
{\large\color{muted}多源遥感表征与联合训练方案\par}
\vspace{13mm}
{\color{teal}\rule{18mm}{1pt}\par}
\vspace{8mm}
{\large\color{ink}作者\enspace\sffamily heyuhang\par}
\vfill
{\color{ink!20}\rule{\linewidth}{.4pt}\par}
\vspace{3mm}
{\small\color{muted}v3.1\hfill 2026-10-27\par}
\end{titlepage}

\pagenumbering{roman}
\hypersetup{pageanchor=true}
\setcounter{tocdepth}{1}
\tableofcontents
\vfill
\note{阅读顺序：第1节了解决策；第2--3节确认数据事实；第4--8节审阅方法与训练；
第9节评测协议，第10节历史训练与评测，第11节100M审计与短程对照结论，第12节生产交付与风险，
第13节讨论投入依据，第14节总结独立目标监督研究及第三、四版模型；附录用于追溯证据和维护文档。
标注为实测样例的质量检查图来自实际数据，其余架构与网格图为原理图。}
\reportpage
\pagenumbering{arabic}

\section{目标、核心决策与适用范围}
\subsection{要解决的问题}
项目目标是把多源年内遥感观测压缩成可复用的密集地理表征，使下游分类、分割等任务可以在冻结嵌入上训练小型预测头。
每个父格输出 \shape{64\times256\times256}，空间步长为5m。向量各维由训练共同形成，没有预先指定的“道路维”或“建筑维”。
设计借鉴 AEF 的多源时空表示和紧凑球面嵌入，但不声称完整复现 AEF，或已经达到其性能\cite{aef}。

\subsection{本版确定的选择}
\noindent\begin{tabularx}{\linewidth}{@{}p{3cm}Y@{}}
\toprule
设计选择 & 技术理由与约束\\\midrule
三分支、一次联合解码 & 低分提供时序语义；多光谱提供局部光谱和空间结构；PAN补充细结构。先各自编码，再对齐融合。\\
固定年度、固定5m输出 & 与现有年度观测和明确产品目标对应，便于建立可验证的第一版合同。\\
端到端训练 & 主干、编码器、年度聚合、联合解码器共同更新；不要求独立训练并永久冻结上采样器。\\
高分可缺失 & 训练模拟删帧、缺模态，推理区分实际高分支持；不以低分复制品填满高分分支。\\
单节点8卡 & 用真实数据测吞吐和显存，再确定骨干容量、观测预算与总训练步数。\\\bottomrule
\end{tabularx}

\subsection{产品边界}
\textbf{5m网格不等于普遍具有真实5m信息分辨率。} 在无同期高分观测的位置，细节由低分输入和学到的先验推断；
多个细结构可能对应近似相同的粗观测，因此无法保证唯一恢复。即使存在高分，云影、几何误差和时间差也可能限制收益。
每份产品同时交付有效性与高分支持信息。

年度嵌入概括一个自然年的观测，不等同于年末快照。年内真实变化不能被一致性损失无条件消除。
全国1\%采样只解决训练与验证数据来源，生产全国连续覆盖仍须准备全国输入。

\note{\status{当前决策}：原4M/100M制品与\texttt{v2.0}标签保持不变。
第三版v3从原100M续训8遍（24,928次更新），加入七项独立目标监督与10m区块加权特征保持；
相对原100M，5m细节$F$平均$+0.54$个百分点（95\%区间$[+0.02,+1.08]$），道路F1不退（$+0.22$），
语义$S$持平（$-0.04$，$[-0.59,+0.43]$）。v3是当前候选版本，尚未发布。
第四版v4以原100M配置加两项实现修复从头预训练200遍，相对原100M细节$F$下降1.92个百分点（$[-3.28,-0.62]$）、
语义$S$下降3.04个百分点（区间跨零），不作为候选。
评测仍基于已访问的弱标签开发集，原S2生产链和独立人工5m真值缺口保留。详见第14节。}
\source{仓库模型与数据实现 [E5]；数据摘要 [E1--E3]；实际训练与升级评测 [E9、E11--E12]；
独立目标资产 [E13]；v3/v4 [E65--E67]。证据编号见附录。}

\reportpage
\section{现有数据与质量状态}
\subsection{核验口径}
当前正式入口为 \texttt{national\_annual5m\_quality\_v6}，覆盖2020、2021两年。
它继承既有低分时序，增加Gaofen 2m PAN并完成吉林一号5m云筛选、全量高分相对配准诊断；不是重新采样的全国网格。
PAN身份根据原始产品名中的PAN标识及日期一致性登记，保留原生640$\times$640数组；不将存储DN宣称为已标定反射率。

\begin{table}[h]
\caption{v6年度正式清单规模。支持表示通过本版保守筛选并存在可读观测，不表示独立科学精度已经验证。}
\centering\small
\begin{tabular}{lrrrrrr}
\toprule 年份 & train & validation & test & 5m支持 & 2m支持 & 三路齐备\\\midrule
2020 & 50,560 & 4,940 & 6,230 & 11,862 & 13,111 & 3,196\\
2021 & 50,560 & 4,939 & 6,230 & 9,279 & 16,694 & 3,082\\\midrule
合计 & 101,120 & 9,879 & 12,460 & 21,141 & 29,805 & 6,278\\\bottomrule
\end{tabular}
\end{table}

\begin{table}[h]
\caption{v6最终保留观测。计数是清单中的观测条目，不等于独立原始卫星景数量。}
\centering\small
\begin{tabular}{lrr}
\toprule 分支 & 来源 & 保留观测数\\\midrule
低分时序 & Sentinel-1 / Sentinel-2 & 1,388,221 / 1,358,421\\
5m多光谱 & JL1GP01 PMS1 / PMS2 & 6,323 / 6,858\\
5m多光谱 & JL1GP02 PMS1 / PMS2 & 8,098 / 4,087\\
2m PAN & GF1B / GF6 / GF1C / GF1D & 1,790 / 28,128 / 1,065 / 2,364\\\bottomrule
\end{tabular}
\end{table}

\subsection{本轮新增处理与隔离}
2m关联目录中142,059条PAN有唯一同景MUX，8卡推理后122,057条通过数值、云影面积和内容对齐筛选；
另有12,056条无法唯一配对，均不进入正式清单。全量跨集合检查按test、validation、train优先级处理，
并对文件摘要、原始景标识和高分像元摘要做重复隔离。

吉林一号四类六波段来源按文件记录的中心波长提取红、绿、近红外，以登记的OmniCloudMask v4运行8卡筛选\cite{ocm}。
102,569条吉林一号观测完成来源标签和原始像元摘要核验；其余高分观测按云影、面积和内容配准规则分别登记。
缺少近红外的五波段来源未纳入本轮清单。v6最终清单保留1,358,421条S2观测；云筛选结果沿用已登记的v3运行证据，不把模型输出当作云影真值。

\begin{figure}[h]
\centering
\begin{tikzpicture}[node distance=7mm]
\node[box,text width=3.6cm] (raw){181,304条\\2m目录记录};
\node[box,text width=3.6cm,right=of raw] (numeric){179,632条\\数值与网格通过};
\node[goldbox,text width=3.6cm,right=of numeric] (selected){122,057条\\PAN筛选通过观测};
\draw[arr](raw)--(numeric);\draw[arr](numeric)--(selected);
\end{tikzpicture}
\caption{PAN处理计数。候选入选不等于云影、辐射与地物配准已通过科学验证。[E6]}
\end{figure}

\subsection{放行边界}
六份清单各抽读100个样本；统计量仅由train内保留观测计算。读取检查不替代云影准确率或地标配准评估。
PAN缺少经验证的云影QA；光学模型对积雪、山地阴影仍需专项复核。

\noindent\begin{tabularx}{\linewidth}{@{}p{4cm}Y@{}}
\toprule 放行标识 & 当前值与意义\\\midrule
\texttt{training\_ready} & true：数据合同、统计量、重复审计和600个真实样本读取已通过；不代表模型已训练。\\
{\scriptsize\texttt{scientific\_quality\_validated}} & false：质量规则和产品语义仍待科学验证。\\
\texttt{annual\_model\_ready} & 数据放行时冻结值为false；此后年度运行时、8卡短跑及两版200轮训练已完成，依据模型证据单独核验。\\\bottomrule
\end{tabularx}
\source{[E1、E3、E8] v6摘要、读取检查与合同门禁；[E2] S2历史运行；核验于2026-09-18。第10节的训练与评测结论依据[E9、E10]，核验于2026-10-02。}

\reportpage
\subsection{本版质量筛选的依据与限制}
\textbf{PAN云影处理。} 全色影像本身缺少红、绿、近红外三通道，不能直接当作OmniCloudMask输入。
本版按父格和原始景标识，关联唯一的同景四波段MUX；按产品惯例读取蓝、绿、红、近红外，
并核对通道、范围、网格、文件摘要及样例\cite{ocm,gfguide}。
142,059组进入8卡处理；12,056条无法唯一配对的PAN隔离。8m云、影和无效区向外扩展一个像元，
按相同地理范围传播至2m网格，再与PAN数值有效掩膜相交。PAN/MUX内容对齐无法确认或偏移超过8m时不用该观测；
保留高分观测须至少20\%预测晴空面积。阈值是本项目的保守工程规则，不是论文提供的精度保证。

\textbf{相对配准诊断。} 在10m诊断网格，将高分与同月S2的有效结构边缘作掩膜归一化互相关；
要求峰值不低于0.35、与非邻近峰的差不低于0.025，且四个局部块至少两个同意整体平移。
256,684条全部登记结果：65,154条可测量且偏移不超过10m，91,589条较大偏移；其余无法可靠判定或重叠不足的观测按最终规则隔离，合计登记为99,939条高分相对配准未解决。
后两类均隔离。该诊断没有独立地面控制点，也不估计局部非刚性畸变；不能据此宣称达到5m绝对或亚像素定位精度。

\begin{figure}[h]
\centering
\includegraphics[width=\linewidth]{/data/heyuhang/processed/national_annual5m_quality_v6/review/report-registration.png}
\caption{实测配准样例R04：同月S2、高分PAN、边缘叠加、诊断平移后的叠加。检测到约51m相对位移。
右图仅用于核查；本版隔离观测，没有自动移动原始影像。[E8]}
\end{figure}

\textbf{云模型的局限。} 已查看42组云影样例及18组配准图。样例中积雪和地形阴影被误屏蔽，
因此不能把模型标签写成云影真值，也不能从这些定向抽例计算准确率。
本版保留保守掩膜，不凭RGB自动恢复雪地；相应的季节、地形和高分覆盖损失须随训练评估报告。
Landsat缺少云影QA，本版不进入训练输入；原文件仍保留。

\begin{figure}[h]
\centering
\includegraphics[width=\linewidth]{/data/heyuhang/processed/national_annual5m_quality_v6/review/report-snow.png}
\caption{实测云筛选样例C31：同景MUX的RGB、云影预测叠加、近红外、PAN。
积雪地形保有纹理却被大范围屏蔽，说明保守筛选会损失可用观测；此例不提供定量误差率。[E8]}
\end{figure}
\source{[E8]全量诊断记录、可追溯图像及目视记录；OmniCloudMask仅提供模型方法，不能替代本区域验证。}

\reportpage
\section{从原始观测到训练样本}
\subsection{数据处理的职责边界}
离线处理负责观测身份、数值含义、质量与索引；模型负责学习特征。保留原始分辨率及实际存储网格，
不把影像提前统一缩放到5m再声称使用原生尺度编码。高分原始归档、候选版本与隔离记录保持可追溯。

\begin{figure}[h]
\centering
\begin{tikzpicture}[node distance=4mm]
\node[box,text width=12.7cm] (a){来源登记：归档摘要、产品、波段、单位、日期与时间支持期};
\node[box,text width=12.7cm,below=of a] (b){像元质量：有限值、NoData、产品QA、云影与饱和检查};
\node[box,text width=12.7cm,below=of b] (c){几何与谱系：CRS、transform、真实内容配准、同景派生关系};
\node[box,text width=12.7cm,below=of c] (d){空间划分：父格与年份绑定；训练区独立统计量；边界缓冲};
\node[box,text width=12.7cm,below=of d] (e){年度manifest：原生数组、逐像元mask、变长高分观测、版本摘要};
\node[goldbox,text width=12.7cm,below=of e] (f){读取与训练验收：真实batch、前反向、8卡短跑、科学质量门禁};
\foreach \src/\dst in {a/b,b/c,c/d,d/e,e/f}\draw[arr](\src)--(\dst);
\end{tikzpicture}
\caption{数据放行与训练验收流程。v6已完成来源至年度manifest和真实读取门禁；年度前反向、8卡短跑与两版长训已有独立运行证据，科学精度验收仍待完成。}
\end{figure}

每条观测至少包含来源ID、父格、日期精度及内容支持期、产品族、波段顺序、单位与缩放、
CRS、transform、有效mask、质量状态、文件摘要和处理版本。只知道月合成的代表日期，不能据此宣称具有逐日观测。
坏包、未知语义与真实缺测分别登记；云影未验证不能写成“晴空”。
当前PAN的有效率仅表示数值有效，不能与S2、吉林一号的预测晴空比例等同。
新版是共享\texttt{data\_root}下的索引型数据集：迁移时须连同清单引用的影像、历史版本掩膜和统计量一起转移。

\subsection{数据放行后的验证任务}
\begin{enumerate}
\item \textbf{产品语义与辐射}：核实各来源波段、单位、scale/offset。统计标准化不能替代辐射标定。
\item \textbf{质量泛化}：S2与吉林一号云筛选需跨区域、季节、地物人工抽验；该模型不能直接套用于单通道PAN。
\item \textbf{绝对定位}：全量高分已完成相对配准诊断并隔离无法确认者；仍需独立地面控制点或权威底图评估绝对精度。
\item \textbf{低分谱系}：高分文件、产品和像元跨集合冲突为零；低分合成影像缺少完整原始景身份，不能宣称全传感器场景级隔离。
\item \textbf{训练闭环}：年度读取器保留PAN原生网格；年度三分支运行时、8卡短跑与两个版本长训已完成，仍须用冻结表征任务验收收益。
\end{enumerate}
\source{[E1] known\_limits与remaining\_model\_work；[E5]年度读取器；[E7--E8]合同、配准与质量记录。}

\reportpage
\section{完整模型：三分支与一次联合解码}
\label{sec:architecture}
主数据流沿用已确定的年度方案。图中省略batch与观测数量维度，采用“通道$\times$高$\times$宽”。
128、32为起步宽度，尚非实测最优参数；输入尺寸以同一1280m父格、规则对齐网格为例。

\begin{figure}[h]
\centering
\begin{tikzpicture}[node distance=5mm]
\node[box,text width=3.8cm,minimum height=12mm] (l0) at (0,0){S2 / S1 年内时序与mask\\Landsat尚未接入};
\node[box,text width=3.8cm,minimum height=12mm] (m0) at (5,0){5m 多光谱观测集合\\$C_5\times256\times256$};
\node[box,text width=3.8cm,minimum height=12mm] (p0) at (10,0){2m PAN 观测集合\\$1\times640\times640$};
\node[box,text width=3.8cm,minimum height=13mm] (l1) at (0,-1.8){分源 Stem\\STP 时空骨干};
\node[box,text width=3.8cm,minimum height=13mm] (m1) at (5,-1.8){5m 专用空间编码器\\输出保持5m网格};
\node[box,text width=3.8cm,minimum height=13mm] (p1) at (10,-1.8){PAN 专用空间编码器\\$32\times640\times640$};
\node[box,text width=3.8cm,minimum height=13mm] (l2) at (0,-3.6){掩膜与日期感知年度聚合\\$F_{10}:128\times128\times128$};
\node[box,text width=3.8cm,minimum height=13mm] (m2) at (5,-3.6){逐景多光谱特征\\$128\times256\times256$};
\node[box,text width=3.8cm,minimum height=13mm] (p2) at (10,-3.6){2m$\rightarrow$5m特征聚合\\通道投影与精化};
\node[box,text width=3.8cm,minimum height=16mm] (l3) at (0,-5.5){插值到5m网格\\$L_5:128\times256\times256$};
\node[box,text width=3.8cm,minimum height=16mm] (m3) at (5,-5.5){日期、质量感知集合聚合\\$G_{5,ms}:128\times256\times256$};
\node[box,text width=3.8cm,minimum height=16mm] (p3) at (10,-5.5){日期、质量感知集合聚合\\$G_{5,pan}:128\times256\times256$};
\node[goldbox,text width=12.5cm,minimum height=16mm] (f) at (5,-7.6){统一5m联合解码器：支持mask、质量控制、通道投影与空间融合\\$H_5:128\times256\times256$};
\node[box,text width=8cm,minimum height=14mm] (v) at (5,-9.5){投影 + vMF瓶颈\\年度嵌入 $64\times256\times256$ + 独立有效性信息};
\foreach \src/\dst in {l0/l1,l1/l2,l2/l3,m0/m1,m1/m2,m2/m3,p0/p1,p1/p2,p2/p3,f/v}\draw[arr](\src)--(\dst);
\draw[arr](l3.south)--(0,-6.65)-|(f.north west);
\draw[arr](m3)--(f);
\draw[arr](p3.south)--(10,-6.65)-|(f.north east);
\end{tikzpicture}
\caption{主架构。高分既可作为可见输入，也可在遮挡后提供训练目标。跨模态只在统一解码器中融合；分支内部多尺度融合属于空间编码。}
\end{figure}

$L_5$是低分语义的网格对齐结果，不是已经读取高分的成品。固定插值没有可训练参数，但允许梯度回传。
端到端训练意味着各可学习模块一起优化，不要求每个几何算子都有参数。不同分支的128维通道不天然语义对齐，须学习融合。
\source{设计依据：AEF/STP\cite{aef}、AnySat\cite{anysat}；三分支合同与一次联合解码为本项目设计，收益待验证。}

\reportpage
\section{高分空间编码与 2m 到 5m 对齐}
\subsection{空间骨干与时序骨干的分工}
高分支线同样需要可训练骨干。卷积可以提取特征；浅层卷积足不足以支撑全国语义表征，则必须实验验证。
主方案将高分建模拆成逐景原生空间编码和带日期的年度聚合，不要求原样复制低分STP。
是否增加高分时空骨干，应依据每父格的有效观测数、季节覆盖和产品一致性决定。

5m编码器的输入输出空间尺寸保持256$\times$256。当前实现使用产品专用输入投影、局部残差空间块与128通道输出。
进一步比较多尺度骨干，利用浅层细节和深层上下文恢复5m网格。FPN提供多层级融合的结构依据，
Swin提供窗口注意力的备选，并不保证任何一个在当前数据上更优\cite{fpn,swin}。

\begin{figure}[h]
\centering
\begin{tikzpicture}[node distance=5mm]
\node[box,text width=3cm] (a){细网格浅层\\边缘、局部纹理};
\node[box,text width=3cm,right=6mm of a] (b){粗网格深层\\更大范围上下文};
\node[goldbox,text width=4.2cm,right=6mm of b] (c){多尺度解码\\结合浅层与深层特征};
\draw[arr](a)--(b);\draw[arr](b)--(c);
\draw[arr](a.south)--++(0,-.5)-|(c.south);
\end{tikzpicture}
\caption{高分编码器的可选内部结构。主图规定外部接口；骨干层数和宽度由同预算消融选择。}
\end{figure}

\subsection{2m PAN 分支的明确接口}
PAN按已验证的单波段产品处理，不复制为RGB。它先在原生网格编码为32$\times$640$\times$640，
再聚合到32$\times$256$\times$256，最后在5m网格精化并扩展至128通道。
2m位置数是5m的6.25倍，所以将大通道扩展放在降采样之后可作为节省计算的起点；实际显存仍需测量。

若采用2m/4m/8m/16m多尺度骨干，本版要求其解码回原生32通道接口再进行5m对齐。
直接从金字塔解码到5m也可试验，但必须登记为独立实现变体。通道扩展和尺寸变化本身不构成语义学习，
特征是否有用取决于训练目标与独立评测。

\subsection{有效掩膜从输入开始生效}
编码前屏蔽无效数值，并将有效性作为条件或使用明确的掩膜卷积策略；编码后传播有效支持。
只在最后乘mask不能保证云影、填充值已不污染邻域。归一化统计也应避免大量填充值支配结果。
不同产品使用自己的辐射处理、统计量和输入投影；主干深层是否共享参数属于实验选择。
\source{原生尺度编码参考AnySat\cite{anysat}；当前实现见[E5] models/annual.py。更深骨干仍需消融。}

\reportpage
\subsection{非整数尺度转换：按地面交叠聚合}
\label{sec:resample}
2m到5m是2.5倍缩小，不能用固定2$\times$2或3$\times$3池化冒充严格对应。
首版使用基于真实transform的面积交叠权重；CRS不同或存在旋转时，须先建立可靠的几何映射。
图\ref{fig:overlap}给出轴对齐情况下的局部示例。

\begin{figure}[h]
\centering
\begin{tikzpicture}[scale=.82]
\fill[teal!15](0,0) rectangle (5,5);
\draw[step=2,gray!70] (0,0) grid (6,6);
\draw[gold,line width=1.6pt](0,0) rectangle(5,5);
\node at (1,1){4};\node at (3,1){4};\node at (4.5,1){2};
\node at (1,3){4};\node at (3,3){4};\node at (4.5,3){2};
\node at (1,4.5){2};\node at (3,4.5){2};\node at (4.5,4.5){1};
\draw[<->] (0,-.5)--node[below]{5m 目标单元}(5,-.5);
\draw[<->] (-.5,0)--node[left]{2m}( -.5,2);
\node[box,text width=6.2cm,anchor=west] at (7,4.3){每个数字为交叠面积，单位m$^2$\\总和为25m$^2$};
\node[box,text width=6.2cm,anchor=west] at (7,2.3){逐特征通道加权汇总\\不是直接平均原始亮度};
\node[goldbox,text width=6.2cm,anchor=west] at (7,.3){无效面积不参与权重\\另存有效覆盖比例};
\end{tikzpicture}
\caption{2m与5m网格的面积交叠示意。其他像元需按自己的交叠计算，不能重复套用这一组权重。}
\label{fig:overlap}
\end{figure}

对目标像元$j$、来源像元$i$、特征通道$c$，定义交叠面积$a_{ji}$、有效支持$m_i$和输入特征$f_{ic}$：
\begin{equation}
 g_{jc}=\frac{\sum_i a_{ji}m_i f_{ic}}{\sum_i a_{ji}m_i},\qquad
 s_j=\frac{\sum_i a_{ji}m_i}{A_j}.
\end{equation}
分母为零时直接返回缺失状态，不能继续相除。$A_j$为目标单元面积，$s_j$是支持比例，
不是模型置信度。分数mask只是有效支持的近似表达，不能宣称知道每个源像元内部无效区域的精确形状。

固定几何权重下该运算对特征可微，梯度能够训练前面的PAN骨干。
面积聚合有低通作用，但不等于理想抗混叠；需检查棋盘格、细线、位移和真实边界。
先编码再聚合有机会保留细结构响应，但无法保证无损保留全部2m信息。

\note{严格区分三个步骤：几何映射确定“在哪里”，空间骨干学习“是什么”，重采样确定“以什么网格表达”。
数组尺寸正确不能代替地物配准，低分回聚一致也不能证明细节正确。}
\source{本节是面积加权聚合的项目定义与几何推导，不归因于某篇神经网络论文；训练效果仍待实测。}

\reportpage
\section{年度集合聚合、融合与球面输出}
\subsection{时间信息在聚合前进入模型}
每景高分先形成5m特征，再结合实际日期、产品身份和质量指标进行掩膜注意力池化。
打分器根据内容与元数据产生权重，无效观测在归一化前排除。单景可以使用；全缺时显式返回零占位与缺失mask，
避免对全掩码序列做softmax产生NaN。

\begin{figure}[h]
\centering
\begin{tikzpicture}
\node[box,text width=3.4cm] (a) at (0,1.2){观测A特征 + 日期A};
\node[box,text width=3.4cm] (b) at (0,0){观测B特征 + 日期B};
\node[box,text width=3.4cm] (c) at (0,-1.2){观测C特征 + 日期C};
\node[goldbox,text width=3.4cm] (pool) at (5,0){质量与支持控制\\集合注意力池化};
\node[box,text width=3.3cm] (out) at (10,0){年度高分特征\\支持mask与观测数};
\draw[arr](a.east)-|(pool.north);\draw[arr](b)--(pool);\draw[arr](c.east)-|(pool.south);\draw[arr](pool)--(out);
\end{tikzpicture}
\caption{年度观测集合聚合。影像、日期与mask一起置换不应改变结果；置换不变不等于忽略时间。}
\end{figure}

年内真实地物变化必须通过日期条件预测及验证表达，而非强制各日期高分特征相同。
年度表征可总结季节轨迹，但注意力结构本身不保证模型真的使用日期；需做日期替换和读出消融。

\subsection{统一解码与缺模态回退}
三路特征形状相同，却没有天然对应的通道语义。将高分特征经过硬支持mask和可学习质量门控后，
与$L_5$拼接，用1$\times$1投影和局部残差卷积得到$H_5$。可保留$L_5$直通路径、使用小残差初始化。
门控不能重新激活硬mask拒绝的像元，也不保证自动拒绝所有坏观测。

高分缺失时，其输入贡献为零，低分空间修正仍可工作，不要求$H_5=L_5$。
本版将有效低分年度上下文作为产品放行条件；低分全缺时即使存在高分，也不自动发布可靠年度嵌入。
高分独立推理须另设数据、训练与验收合同。

\subsection{最终64维输出}
通道投影与L2归一化统一放在vMF瓶颈；各分支不新增独立的L2瓶颈。训练用解码器读取该64维输出。
现有实现使用环境空间高斯扰动后重新归一化，推理不加噪声，不能写成精确vMF采样。
共享的\texttt{log\_kappa}也不是逐像元校准置信度。单位范数不能代替有效mask\cite{vmf}。
\source{当前实现见[E5] bottleneck.py、annual.py；vMF噪声不是独立置信度校准。}

\reportpage
\section{训练信号与防止信息泄漏}
\subsection{没有天然的5m嵌入真值}
首版依赖真实观测自监督，不要求现成的5m教师特征或人工标签。训练时从年度观测中留出目标，
其余观测产生嵌入，再由日期、产品条件头预测目标。监督信息来自影像本身，语义是否有用由独立下游任务判断。
AEF提供来源重建的参考；TESSERA与Barlow Twins提供观测子集和冗余约束的参考\cite{aef,tessera,barlow}。

\begin{figure}[h]
\centering
\begin{tikzpicture}
\node[box,text width=3.2cm] (input) at (0,0){年内观测集合\\先留出目标};
\node[box,text width=3.2cm] (model) at (4.5,0){可见输入\\三分支模型};
\node[box,text width=3.2cm] (embed) at (9,0){64维年度嵌入\\唯一信息瓶颈};
\node[goldbox,text width=4cm] (target) at (0,-2){被留出观测\\不进入输入、跳连与统计};
\node[box,text width=4cm] (decode) at (9,-2){日期/产品条件读出\\预测低分及高分目标};
\node[labeltext] (loss) at (4.5,-2){有效目标上的损失};
\draw[arr](input)--(model);\draw[arr](model)--(embed);\draw[arr](input)--(target);\draw[arr](embed)--(decode);
\draw[arr](target)--(loss);\draw[arr](decode)--(loss);
\end{tikzpicture}
\caption{训练信息流。空间编码器的内部跳连允许存在，但目标读出不得绕过64维瓶颈复制影像。}
\end{figure}

\noindent\begin{tabularx}{\linewidth}{@{}p{2.7cm}YY@{}}
\toprule 目标 & 监督与作用 & 限制\\\midrule
低分遮挡预测 & 原生网格上的有效观测；保留光谱和时序信息 & 不能把放大的Landsat像元当独立高分真值\\
高分遮挡预测 & 5m光谱；PAN聚合至5m的亮度或结构统计 & 不要求5m瓶颈无损还原2m全图\\
子集一致性 & 同父格、同年度两个观测子集的公共语义 & 缺高分视图不逐5m像元强制细节相同\\
冗余与坍塌控制 & 监测方差、相关性、有效秩；训练投影头 & 跨父格统计，相邻像元不能冒充独立样本\\
可选辅助标签 & 有可靠时间和覆盖的地物标签 & 未知不作负例；与独立评测严格隔离\\\bottomrule
\end{tabularx}

损失权重、结构目标和抽样比例均为待标定配置。按有效目标数、产品和来源归一化，避免观测多的来源支配优化。
自动边缘可能来自阴影或噪声，不能当作道路真值。训练统计量只来自训练区，不能重复使用目标内容拟合条件。

\subsection{遮挡与谱系检查}
空间遮挡在编码前、按地理位置映射到相关分支；同景别名、重采样派生及相关产品需登记，防止绕过留出规则。
保留另一模态作为上下文时明确标为跨模态预测。实际日期不随机伪造；发生真实变化的区域不施加强细节一致性。

\reportpage
\section{训练顺序、8卡资源与可选升级}
\subsection{从数据到可发布模型}
\noindent\begin{tabularx}{\linewidth}{@{}p{2.5cm}YY@{}}
\toprule 阶段 & 主要产物 & 进入下一阶段的依据\\\midrule
G0 数据放行 & 产品/质量/配准与谱系审计、版本化manifest & 接口和科学质量证据分别齐备\\
G1 模型闭环 & 年度接口、三分支前向与反向、缺模态处理 & 活跃分支有有限梯度，目标无泄漏\\
G2 8卡短跑 & 吞吐、显存、有效样本与目标计数 & 真实I/O和缺模态DDP稳定，能够恢复\\
G3 效果对照 & 固定空间留出、种子和下游头的消融结果 & 高分收益可重复、成本可接受\\
G4 正式训练 & 固定结构与预算、训练划分长训 & 验证通过；最终test仅报告一次\\
G5 全国生产 & 固定权重及输入版本的逐父格输出 & 几何、数值、支持状态和接缝验收\\\bottomrule
\end{tabularx}

可以短暂预热新增模块，再联合训练，但不设置“先练好并永久冻结上采样器”的硬依赖。
已有权重仅在波段、预处理、时间、参数形状与来源一致时用于初始化；形状相同不等于语义兼容。
全国1\%样本的validation/test始终留出，不因扩大训练而合并回训练集。

\subsection{算力使用与可复现性}
统一单节点8卡。按产品与尺寸分桶，限制单样本高分观测数，并跨轮次覆盖不同季节；
生产和验证使用确定性选帧。采用混合精度、梯度累积和必要的激活检查点，容量依据实测调整。
报告GPU小时、样本曝光、有效目标数与I/O等待，相同步数不代表相同计算量。

可缓存原始数组、mask和几何权重，但不可固定可学习特征缓存后仍称骨干端到端更新。
缺模态batch需检查DDP未使用参数处理，跨卡相关性统计按真实优化器步数更新。
源码工作区存在未提交修改，checkpoint来源不能只写Git HEAD，须保存实际源码摘要。

\subsection{升级只按单项实验进入主线}
\noindent\begin{tabularx}{\linewidth}{@{}p{3cm}Y@{}}
\toprule 候选方法 & 用途与边界\\\midrule
DySample\cite{dysample} & 学习采样偏移，比较是否优于固定插值；偏移不改变产品地理网格，也不替代配准。\\
FeatUp\cite{featup} & 借鉴高分引导与多视图一致性；跨传感器日期差不等同于同图增强，不采用全国逐图优化作为首版。\\
I-JEPA / 冻结教师\cite{ijepa} & 提供潜在目标；冻结分支使用已登记checkpoint，EMA作为独立对照并保存其更新状态。\\
多尺度骨干 & 比较上下文与细节收益；不同时改动所有分支、损失和采样策略。\\\bottomrule
\end{tabularx}
\source{本节资源分配为实施设计；两版训练的实测吞吐与显存见第10--11节，不能用短跑吞吐直接承诺完整重训耗时。}

\subsection{年度训练目标与融合约束}
三分支首先使用留出观测重建作为主自监督目标；教师--学生潜变量预测是可选第二阶段：
学生只看到随机留出的观测子集，教师由已验收并冻结的 M0/M1 checkpoint 产生目标嵌入；教师输入包含另一观测子集，
且两者在空间、月份和来源上严格隔离。这样目标不必经过像素级配准到5m，也不会要求模型把云影或时间变化复制回输入。
AnySat采用JEPA和尺度自适应空间编码器处理异构分辨率，结果支持将“预测潜变量”作为主自监督方向，
而不是把所有高分细节都压成重建误差\cite{anysat}。

高分融合采用“残差 + 门控 + 低分锚点”。先得到只由低分时序产生的$L_5$，再由高分分支预测残差$R_5$和
门控$\alpha_5$：$H_5=L_5+\alpha_5R_5$。$\alpha_5$由有效支持、日期距离、配准质量和内容一致性共同约束，
并以小值初始化；高分缺失时严格退化为$L_5$。这样模型不会因为某个高分观测存在就覆盖稳定的低分语义。

年度高分聚合同时计算加权特征、有效覆盖比例、观测数和日期差，
将这些量作为显式条件输入；训练时随机删除观测、随机保留月份，并加入重复观测不变性约束。TESSERA采用稀疏随机时间采样
和Barlow Twins约束来减少对具体观测子集的依赖，这一思想适合本项目的缺云和高分稀疏场景\cite{tessera,barlow}。

2m到5m的几何聚合和细节增强解耦。面积交叠权重负责地理守恒；高分编码器学习局部结构；可学习采样器
（如DySample）作为独立消融，不能改变输出网格或修正未解决的配准误差\cite{dysample}。主模型使用固定几何聚合。

\begin{table}[h]
\caption{算法消融顺序。每次只改变一项，所有组使用相同空间划分和训练预算。}
\centering\small
\begin{tabularx}{\linewidth}{@{}p{1.2cm}p{4cm}Y@{}}
\toprule 组别 & 改动 & 主要判断\\\midrule
M0 & 当前重建目标 + 固定几何聚合 & 建立可复现基线，记录缺高分区域退化\\
M1 & M0 + 低分锚点残差门控 & 检查高分是否改善细节且不破坏低分语义\\
M2 & M1 + 观测子集教师--学生潜变量预测 & 检查跨日期、跨来源和缺云稳定性\\
M3 & M2 + 覆盖/日期/配准质量条件 & 检查模型是否学习质量差异，而不是盲目平均\\
M4 & M3 + 可学习采样器或多尺度骨干 & 只有在M3稳定领先时才承担额外显存和解释成本\\\bottomrule
\end{tabularx}
\end{table}

M1构成首个正式训练主线，M2/M3在教师 checkpoint 和质量条件验收后进入对照，M4作为后续实验。每组至少报告四类指标：低分-only、5m-only、PAN-only和完整输入；
同时报告有高分与无高分、云覆盖高与低、相对配准误差小与大的空间块。若高分只提高重建损失而不能提高冻结嵌入下游结果，
不进入正式模型。若M2导致表示坍塌或跨年度不稳定，则退回M1并增加来源归一化和目标遮挡比例消融。

\subsection{方法边界}
2m PAN不作为“真实5m语义真值”监督：PAN缺少光谱，且高分与低分存在时间差和残余几何误差。
已完成主线及本轮ABC未将外部AlphaEarth嵌入用作逐像元监督，目前只用其作外部冻结表征对照；
未来若增加外部教师，须另登记目标来源并作独立消融。
AlphaEarth公开描述的是年度、多源输入的64维嵌入场，说明年度压缩和条件读出的产品方向有参考价值，但不能替代本地训练证据\cite{alphaearth}。

\reportpage
\section{验证方法、对照与停止条件}
\subsection{验证什么，而不是仅看训练损失}
\noindent\begin{tabularx}{\linewidth}{@{}p{1.2cm}YY@{}}
\toprule 组别 & 模型/输入 & 要回答的问题\\\midrule
A0 & 低分时序 + 固定5m网格解码 & 没有高分时的基础能力与成本\\
A1 & A0 + 5m多光谱 & 光谱和高分结构是否带来收益\\
A2 & A0 + 2m PAN & 全色细结构是否带来收益\\
A3 & 完整三分支 & 两类高分是否互补\\
B组 & 在优选模型上逐项改变骨干、损失或上采样 & 收益能否归因，开销是否值得\\\bottomrule
\end{tabularx}

各组使用相同的空间划分、标签、预测头和调参预算。至少两个预训练随机种子，报告各次结果而非只选最好一次。
比较冻结嵌入上的地表覆盖、建筑、道路、水体等任务；可用性不足的标签先审计，不默认已经具备完整标注。
采用任务适合的IoU、F1等指标，以及真实坐标边界评估。置信区间按空间块抽样，不能把相关像元当独立样本。

高分收益至少区分：训练时见过高分的区域、空间独立区域、推理有高分与无高分。
可以加入同年度AEF公开产品作为外部基线，但须先确认产品实际可用、标签与网格一致，不能把参考论文数字搬来作本地对照。

\subsection{必须通过的反事实与接口检查}
\begin{enumerate}
\item 高分全缺、单源缺失、全部输入无效时，输出状态正确且无NaN；无效区域不能被当作有效产品。
\item 同时置换影像、日期与mask，年度集合输出在数值容差内不变；重复观测不得重复加权。
\item 修改被屏蔽像元和留出目标，不应通过未授权路径影响输入编码；同景派生文件同样受约束。
\item 以已知位移和独立地标检查几何，检验父格边缘接缝；不得用允许任意平移的损失掩盖位置误差。
\item 检查活跃分支梯度、门控饱和、有效秩、低分任务退化；与仅插值对照比较。
\end{enumerate}

\subsection{放行与停止}
在效果实验前登记主指标、最低样本量、可接受低分退化和实际成本上限。当前没有证据支持写死“提升若干百分点”，
也没有依据承诺训练时长。若高分只降低重建损失、未改善冻结表征，则不据此放行；
若高分效果依赖配准错误、特定传感器或空间泄漏，先修数据与协议。
缺少合格独立标签时，结论只能停留在工程可运行，不能宣称正式模型优于基线。
\source{本节规定验收边界；第10节历史弱标签评测与第11节升级实验采用各自登记的具体协议，不能合并不同探针口径的数字。}

\section{4M历史训练、冻结嵌入评测与外部对照}
\subsection{正式训练与收敛}
\noindent\begin{tabularx}{\linewidth}{@{}p{3.4cm}Y@{}}
\toprule 项 & 实测值\\\midrule
模型与规模 & \texttt{annual5m\_v1}，4{,}010{,}057参数（推理路径3{,}636{,}257）；输出64维、5m、$64\times256\times256$\\
硬件与耗时 & 单节点8卡A100-SXM4-40GB；200 epoch，337{,}753\,s（约3.91天）\\
样本与步数 & 20{,}224{,}000样本曝光，632{,}000次优化器更新，全局吞吐59.9样本/s\\
数据等待 & 25{,}864\,s，占wall time的7.7\%\\
数值健康 & 非有限梯度步0次；峰值显存23.3\,GB/卡；全程无异常退出\\
来源 & 配置SHA \texttt{04e1fe06}$\cdots$，Git SHA \texttt{9fb2e67}（标签\texttt{annual5m-v1.0}）\\
掩蔽重建误差 & 第1轮0.682$\to$第10轮0.515$\to$第200轮0.466；批内标准差0.249$\to$0.108\\\bottomrule
\end{tabularx}

\source{[E9] 训练日志与checkpoint。重建误差为训练批次上的逐轮均值，尚未在验证集与测试集上单独计算，不能据此判断过拟合。}

\subsection{评测口径}
本轮评测使用外部弱标签，冻结上游权重，只训练参数极少的探针；这与“数据合同通过不等于模型已验收”的
既有边界一致，本节结论不得外推为独立精度认证。

\begin{enumerate}
\item \textbf{标签来源与同源风险}：地表覆盖与水体取自ESA WorldCover 2021，水体另有JRC交叉来源；建筑、道路取自
OSM（Geofabrik 2022-01-01快照），按5m栅格化。WorldCover由Sentinel-1/2训练得到，与本项目输入同源，
其上的绝对分数偏乐观，只用于方法间相对比较。
\item \textbf{探针与划分}：逐像元线性回归与$k$近邻（$k=5$）用于点协议；线性与$3\times3$卷积头用于5m稠密协议。
训练集内评测按patch五折交叉验证；与外部基线的对比按城市分组划分，使测试城市对探针完全未出现。
\item \textbf{统计}：置信区间按patch做1000次聚类bootstrap；模型间比较使用配对bootstrap，
报告差值的95\%区间。同一patch内像元相关，未按像元计算独立性。
\item \textbf{分辨率对齐}：与10m标签对比时，先对嵌入做$2\times2$平均池化；5m评测直接使用原生输出，
并把外部10m基线双线性上采样到5m网格，同时以最近邻上采样作为敏感性对照。
\end{enumerate}

\subsection{与外部年度嵌入基线的对照}
外部基线为公开的AlphaEarth年度嵌入（Google DeepMind，64维、10m、2021年）。两套嵌入在同一批patch、
同一折划分、同一探针下评测。需说明：外部基线的训练数据很可能已覆盖本项目的测试区域，因此“未见区域”一组
对外部基线并非严格未见，该对照只能说明方法之间的差距，不能据此判定泛化能力。

\begin{table}[h]
\caption{历史10m点协议，按patch五折；少样本行为三个种子均值，全量行为单次折外结果及其配对95\%区间。}
\centering\small
\begin{tabularx}{\linewidth}{@{}p{2.6cm}YYYY@{}}
\toprule 任务与指标 & 本项目 & 外部基线 & 原始波段 & 配对差值（本项目$-$外部）\\\midrule
水体F1（每类全量） & 0.923 & 0.932 & 0.817 & $-0.009$ [$-0.036$, $+0.022$]\\
地表覆盖mIoU（每类10） & 0.344 & 0.411 & 0.297 & $-0.067$（均值差）\\
地表覆盖mIoU（每类100） & 0.395 & 0.474 & 0.393 & $-0.079$（均值差）\\
地表覆盖mIoU（每类全量） & 0.481 & 0.533 & 0.428 & $-0.052$ [$-0.082$, $-0.024$]\\\bottomrule
\end{tabularx}
\end{table}

历史脚本的少样本区间仅使用seed 0，而表中主分数为三种子均值。旧版报告曾把二者并列，
导致差值与主分数不相减。本版改为报告可核对的均值差，不给出未正确汇总三种子的区间；
原始seed 0统计保留在[E10]。全量行不存在该混用。

差距集中在依赖物候的类别（耕地IoU 0.660对0.836、草地0.637对0.745、建筑0.647对0.782），
水体、裸地点分数接近。该结果提示应优先诊断物候相关类别，但不能由外部对照分数单独
归因于时间上下文、多年信息或某个结构模块。

\source{[E10] 110个训练集patch上的折外探针结果与配对检验。}

\begin{figure}[p]
\centering
\includegraphics[width=\linewidth]{/data/heyuhang/xuannv_embedding/outputs/eval/quick_rapid124/figures/water.png}
\caption{实测样例：水体提取（10m，折外线性探针）。每行一个训练集patch，列依次为S2年度中位数真彩色、
WorldCover水体标签、本项目嵌入预测、本项目误差、原始波段基线误差、嵌入前三主成分。误差图中蓝色为正确水体，
橙色为误报，品红为漏检。样例按水体面积分位选取，不按效果挑选。}
\end{figure}

\begin{figure}[p]
\centering
\includegraphics[width=\linewidth]{/data/heyuhang/xuannv_embedding/outputs/eval/quick_rapid124/figures/landcover.png}
\caption{实测样例：地表覆盖分类（10m，WorldCover 11类，折外线性探针）。列依次为真彩色、标签、本项目嵌入预测、
原始波段基线预测、嵌入前三主成分。样例取类别最丰富的五个patch，单块精度（0.44--0.86）低于总体的0.84。
第3行工业厂房在标签中为“裸地”，属标签误差；第1、5行湿地与灌木未被识别，与该两类IoU接近0一致。
多张主成分图的右缘与下缘可见细色带，提示patch边界伪影，需在拼接验收中单独检查。}
\end{figure}

\subsection{5m原生协议与外部基线的对照}
5m网格使用OSM道路、水体与建筑标签。外部基线在10m上采样到5m网格；同时给出本项目自身池化到10m再上采样
的结果（下表“10m对照”），用于把“分辨率”与“模型内容”两种贡献分开。

\begin{table}[h]
\caption{5m稠密协议，按城市分组五折；$3\times3$卷积头。道路186个patch，建筑38个（标签覆盖不足），水体80个。}
\centering\small
\begin{tabularx}{\linewidth}{@{}p{3.4cm}YYYY@{}}
\toprule 任务与分组 & 本项目5m & 10m对照 & 外部基线 & 配对差值（本项目$-$外部）\\\midrule
道路F1（全部） & 0.228 & 0.230 & 0.204 & $+0.024$ [$+0.008$, $+0.038$]\\
道路边界F1（全部） & 0.321 & 0.323 & 0.270 & $+0.051$ [$+0.035$, $+0.069$]\\
道路F1（未见城市） & 0.213 & 0.218 & 0.170 & $+0.043$ [$+0.012$, $+0.077$]\\
水体F1（未见城市） & 0.710 & 0.701 & 0.570 & $+0.140$ [$+0.043$, $+0.235$]\\
建筑F1（全部） & 0.251 & 0.248 & 0.259 & $-0.008$ [$-0.066$, $+0.040$]\\\bottomrule
\end{tabularx}
\end{table}

道路与城市水体上的领先在统计上成立，但同一张表显示：本项目5m与自身的10m对照几乎一致，
提示该领先主要与嵌入内容有关，尚缺少稳定的分辨率增益证据。融合过程的逐层测量提供了进一步诊断线索。

\begin{table}[h]
\caption{高分辨率信息在各融合阶段的保留情况（16个含高分观测的patch均值）。}
\centering\small
\begin{tabularx}{\linewidth}{@{}p{4.4cm}YY@{}}
\toprule 阶段 & 特征范数 & 5m频段能量占比\\\midrule
低分主干$L_5$（10m上采样） & 5.30 & 2.6\%\\
高分全色分支（2m输入） & —— & 18.5\%\\
高分残差$\alpha_5R_5$（门控均值0.395） & 0.68 & 8.3\%\\
最终嵌入（vMF球面单位向量） & 1.00 & 3.1\%\\\bottomrule
\end{tabularx}
\end{table}

高分分支确实编码了细节，但残差幅值仅为低分主干的约13\%，叠加后被主干淹没：含高分patch的最终高频占比
（3.0\%）与不含高分的（3.2\%）没有区别。该历史4M产品可概括为“10m质量的语义、5m格式的输出”，
尚未兑现高分辨率相对外部基线唯一的结构性差异。

\source{[E9、E10] 原生5m导出、OSM标签（含860MB全国快照）与逐层钩子测量；城市标签共3901栋建筑、
8092条道路、600个水体面，国内OSM建筑覆盖偏稀，建筑一组不作为结论依据。}

\subsection{已知问题与后续动作}
\begin{enumerate}
\item \textbf{高分融合稀释细节（首要）}：按M1方向改造——融合前统一各分支尺度、把高分特征拼接进5m解码器
而非仅作小残差、提高高分重建权重并对含高分样本过采样。用“本项目5m相对自身10m的增益”和“含高分时的高频能量”
作为验收量；短程继续训练只作筛查，不能据此保证从头长训的最终收益。
\item \textbf{土地覆盖落后于外部基线}：优先消融时间建模，包括月合成是否损失物候、是否需要多年上下文、
随机遮蔽比例是否过高。
\item \textbf{5m标签不足}：国内OSM在多数地点建筑稀疏，建筑与道路的绝对精度不可比；需要可靠5m真值
（高分影像人工精标或高质量公开轮廓）才能独立证明分辨率收益。
\item \textbf{推理输入规格未闭环}：训练低分数据的波段定义未随数据发布（\texttt{band\_semantics}标注为待核实），
现有导出命令也只能按训练清单导出，尚不能指向任意manifest；两项未解决前不具备生产推理条件。
\item \textbf{基线对照的公平性}：两者训练地域、观测类型、目标与规模不同，样本或观测计数不能直接互换；
外部基线训练集还可能覆盖本地测试区域。结论应表述为“在同一评测协议下与外部基线的差距”，
不得表述为全面优劣。
\end{enumerate}

\reportpage
\section{100M升级：实现审计与受控实验}
\note{本节压缩记录2026-10-15的100M审计与短程对照；完整过程、成本、逐组头种子复核与诊断保留在归档标签
\path{archive/annual-research-20261019}及[E12]中。涉及旧多PAN实现的多卡对照缺少逐rank参数身份保证，
只保留为制品分数，不作因果归因（见第14.1节）。}

\subsection{两个已训练版本}
\noindent\begin{tabularx}{\linewidth}{@{}p{4cm}YY@{}}
\toprule 项 & 原4M & 原100M\\\midrule
总参数 & 4{,}010{,}057 & 104{,}216{,}905\\
STP空间/时间宽度与层数 & 128 / 128 / 2 & 1280 / 512 / 4\\
时空STP参数 & 3{,}428{,}480 & 103{,}635{,}328\\
高分编码器 / 联合解码参数 & 107{,}392 / 67{,}328 & 107{,}392 / 67{,}328\\
训练轮数 / 优化器更新 / batch & 200 / 632{,}000 / 32 & 200 / 632{,}000 / 32\\
耗时 / 吞吐 & 3.91天 / 59.88样本/s & 4.72天 / 49.64样本/s\\\bottomrule
\end{tabularx}

容量增加几乎全部进入低分时空主干，高分编码与5m解码没有扩大，所以不能把总参数的26倍理解为5m结构容量的26倍。
两版同时改变了STP深度、宽度和重采样顺序，不是严格的规模消融。训练记录101{,}120条，对应50{,}562个唯一位置；
有任意高分的样本占35.82\%，两类高分同时具备的只有5.02\%。
\source{[E9、E11] 两版checkpoint、配置与日志；[E1]训练清单分支支持计数。}

\subsection{两项实现缺陷}
\textbf{遮挡区域的上下文梯度（A）。} 历史池化用人工遮挡mask把STP在遮挡块内推断出的特征再次置零，
块内重建对STP和瓶颈投影没有梯度。修复把“编码器可见性”与“真实有效性”分开：
目标值仍在编码器前隐藏，STP之后按真实有效性聚合（\path{annual_context_pooling}）。
真实权重诊断确认旧路径梯度为0、新路径非零；这只证明梯度路径打通，不证明下游收益（第14.3节的从头训练结果为负）。

\textbf{PAN可见性投影误删。} 10m可见性投影到2m时，Float32面积权重的行和误差可达$4.58\times10^{-5}$，
旧代码以$1-10^{-6}$判断完全可见，导致48.56\%的有效PAN像元被误删，生产导出同样受影响。
修复用整数网格包含关系或Float64权重传递mask（\path{annual_exact_mask_coverage}），误删降为0。

\begin{figure}[ht]
\centering
\includegraphics[width=.9\linewidth]{/data/heyuhang/xuannv_embedding/outputs/annual_abc/20261015_exact/visualization_inputs/real_pan_input_report.pdf}
\caption{预先按哈希选定的开发位置：原生PAN、历史可见输入、修复后可见输入与OSM标签。
白色条纹为误删输入。OSM明显漏标建筑，不能把未标注区当可靠负样本。本图不构成精度评测。}
\end{figure}
\source{[E12]真实梯度、原生PAN诊断与源码快照。}

\subsection{短程配方对照与选择}
从原100M出发，用相同8{,}192条年度记录、2{,}048次更新（约原训练的0.3\%）比较十一组配方：
同预算control、A、AB（高分可见/留出目标，可见权重$\lambda=0/0.1/0.3$）、
ABC（再加冻结4M教师、EMA视图一致性或两者），以及两组高分权重3的追加。
评测与第14节相同：124个语义位置、203个5m城市位置，城市/空间块配对bootstrap。

\noindent\begin{tabularx}{\linewidth}{@{}p{2.6cm}YYYY@{}}
\toprule 组别 & 10m水体F1 & 覆盖mIoU & 5m道路F1 & 建筑F1\\\midrule
4M参考 & 0.9064 & 0.4487 & 0.2182 & 0.2562\\
原100M & 0.8538 & 0.4315 & 0.2308 & 0.2390\\
仅修PAN输入 & 0.8611 & 0.4378 & 0.2302 & 0.2487\\
control & 0.8551 & 0.4372 & 0.2330 & 0.2423\\
A & 0.8666 & 0.4788 & 0.2242 & 0.2388\\
AB，$\lambda=0$ & 0.8785 & 0.4912 & 0.2192 & 0.2422\\
ABC，联合 & 0.8652 & 0.4886 & 0.2339 & 0.2420\\\bottomrule
\end{tabularx}

\begin{itemize}
\item 学习型改动普遍提高语义，但道路F1下降；AB各组道路差的区间完全低于0。
覆盖mIoU的增量约九成来自只有两个patch的雪冰类，不能直接解读为通用语义提升。
\item 没有候选同时通过语义与细节约束。按预先锁定的规则选定\texttt{original\_exact}：
原100M参数不变，只修PAN输入；相对原100M，只有覆盖mIoU的区间大于0（$[+0.0012,+0.0121]$）。
\item 种子17复验中，AB道路仍下降，ABC联合组的细节差值方向随种子翻转，不构成稳定双提升。
\item 80个封存位置的空间确认中，各差值区间均跨0；原4M在该确认集上道路F1比原100M高3.4个百分点。
\end{itemize}
这些短程结论只约束该初始化下的继续训练配方，不能替代从头训练。它直接引出了第14节的独立目标监督研究。

\begin{figure}[ht]
\centering
\includegraphics[width=.9\linewidth]{/data/heyuhang/xuannv_embedding/outputs/annual_abc/20261015_exact/figures/semantic_fine_tradeoff.pdf}
\caption{全部十一组训练及参考制品在开发协议下的语义/细节点估计。B表示AB，T冻结教师，V为EMA，HR3为高分权重3；4M仅为参考。}
\end{figure}
\source{[E12] 冻结plan、选择规则、\texttt{selection.json}、种子17复验与空间确认结果。}

\reportpage
\section{实现边界、生产交付与风险}
\subsection{代码实现与验证边界}
\noindent\begin{tabularx}{\linewidth}{@{}p{2.6cm}YY@{}}
\toprule 模块 & 已有代码基础 & 年度状态与验收余项\\\midrule
低分模型 & STP、月度聚合、插值卷积头 & 年度STP汇总与128通道接口已实现\\
高分编码 & 原生网格卷积，支持先编码后缩放 & 5m/PAN产品适配、mask传播与面积聚合已实现\\
观测融合 & 按目标月聚合和残差融合 & 日期/质量年度聚合、两类高分联合解码已实现\\
球面瓶颈 & 64维投影、归一化、训练扰动 & 接入年度路径与输出有效性\\
年度数据 & 原生网格读取、2m接入、5m云筛选、全量高分相对配准与最终统计量 & 传感器语义复核、独立绝对定位评估\\
训练与导出 & 严格配置/checkpoint、DDP、月度导出 & 年度监督、生产导出及两版长训已完成；升级收益与科学精度分别验收\\\bottomrule
\end{tabularx}
\source{[E5]是实际源码摘要；上表描述代码能力；实际运行范围见[E9、E11--E12]，不由代码存在推断科学精度。}

本轮软件检查在\texttt{58f4374}对应源码上完成：393项测试通过、2项跳过，Black、Ruff与
仓库策略检查通过；wheel独立安装路径核验、11组CLI帮助及CPU单步前向、反向、保存通过。
安装验证共享既有依赖，不是全新依赖环境认证。跳过项涉及未取得的真实431键旧制品和
未安装的OmniCloudMask依赖，不能登记为通过；本轮也不替代NPU平台验收。
真实栅格的8卡CUDA恢复检查已完成：从实际EMA周期checkpoint继续两次更新、64次样本曝光，
Adam最大步数与EMA计数均由512前进到514，426个EMA状态键与学生完整匹配；
模型、损失模块与EMA状态张量均有限，学生及EMA投影权重均有实际变化，非有限梯度步为0。
原checkpoint摘要保持不变，配置、来源与区域身份一致。该检查验证状态可恢复并继续更新；
当前checkpoint不记录RNG状态，因此不宣称逐位复现随机轨迹，两步后的epoch标记也不代表完整一轮。
\source{[E12] \texttt{quality/summary.json}、完整测试日志、wheel验证及
\texttt{ema\_cuda\_resume/result.json}与真实恢复日志。}

\subsection{全国生产合同}
每个父格和年份保存64$\times$256$\times$256嵌入，以及有效mask、5m多光谱支持、PAN支持、
有效覆盖与观测数量。记录年份、CRS、transform、波段/来源注册表、时间规则、数据版本、统计量、配置与模型摘要。
推理去掉重建头与一致性头，固定参数和统计量，关闭瓶颈噪声，不按地区继续训练。

FP16单父格裸嵌入容量由64$\times$256$\times$256$\times$2字节计算为8MiB，另计掩膜、元数据及容器开销。
全国容量须依据实际父格数测算，不把训练样本数当全国网格数。按块写出、校验后原子发布，失败可续跑；
重叠区域和无效边缘采用固定规则。模型改变后必须区分版本，不能拼接未经对齐的不同嵌入空间。
即使参数不变，PAN输入可见性修复也会改变输出；相应导出须登记新的配置与输入策略版本，
不能与历史输入策略产生的嵌入直接混拼。

\subsection{主要风险及应对}
\noindent\begin{tabularx}{\linewidth}{@{}p{3cm}Y@{}}
\toprule 风险 & 处理与验证\\\midrule
高分覆盖偏向某些地区 & 分层抽样与空间留出；分别评估推理有/无高分。\\
纹理捷径与云影污染 & 有效QA、目标留出、独立语义评估；不能靠视觉锐利判断好坏。\\
年度变化被平均抹除 & 日期条件预测、季节分桶；避免跨日期无条件一致性。\\
多版本数据含义冲突 & 每条统计标明数据版本与证据；不同流水线不合并计数。\\
复杂模块增加维护成本 & 先固定三分支基线；每次升级都保留可归因的消融与回退。\\\bottomrule
\end{tabularx}

\reportpage
\section{参考工作与本项目的采用边界}
本报告的结构组合属于项目设计。参考论文提供方法依据，不等于相关模型、预训练权重或下游收益已经在本地复现。

\noindent\begin{tabularx}{\linewidth}{@{}p{2.6cm}YY@{}}
\toprule 工作 & 原方法要点 & 本项目采用位置及未直接采用部分\\\midrule
AEF\cite{aef} & 多源时空表征，STP与紧凑球面嵌入；来源重建等目标 & 低分时序主干与64维输出的依据。三分支5m融合是本项目设计。\\
AnySat\cite{anysat} & 尺度自适应空间编码器与JEPA，处理异构观测 & 原生尺度、产品投影和地面范围对齐。首版不声称复现其JEPA全目标。\\
TESSERA\cite{tessera} & S1/S2时间序列，独立稀疏观测子集与修改的Barlow目标 & 观测子集一致性；扩展到高分缺失是项目设计，不是其空间上采样方法。\\
Barlow Twins\cite{barlow} & 两视图交叉相关的对角一致与非对角冗余约束 & 训练投影头的候选正则；需处理像元相关与跨卡统计。\\
FPN / Swin\cite{fpn,swin} & 多尺度横向连接；分层移位窗口注意力 & 高分空间骨干候选。PAN/多光谱不能直接等同RGB预训练输入。\\
Scale-MAE\cite{scalemae} & GSD感知位置编码与多尺度/频率重建 & 保留物理尺度与结构目标的参考；未实现全方法前不称为使用该模型。\\
Galileo\cite{galileo} & 全局/局部目标联合，区分语义与细结构的学习尺度 & 本轮C只采用粗尺度目标的消融思路，未复现其PatchDisc和完整目标。\\
OlmoEarth\cite{olmoearth} & 冻结随机投影目标、模态感知遮挡与潜在空间判别 & 提醒重建损失不等于表征质量；本轮未实现其完整Latent MIM Lite。\\
DINOv3\cite{dinov3} & 用Gram anchoring保护长训中的密集特征关系 & 为细节保持目标提供候选依据；本轮未实现，也未使用其权重。\\
SALT\cite{salt} & 先训练目标编码器，再冻结教师进行潜变量预测 & 冻结教师分支采用两阶段接口；EMA属于本轮另一独立消融。\\
Presto\cite{presto} & 面向像元时序的轻量 Transformer 与结构化掩膜 & 作为低成本非 STP 对照，不作为正式5m主干。\\
vMF研究\cite{vmf} & 监督学习的球面几何与分布建模 & 瓶颈的数学背景；当前高斯投影实现不等于精确采样或校准不确定性。\\\bottomrule
\end{tabularx}

\subsection{方法选择由产品需求与数据条件驱动}
固定年度5m产品需要同时表达年内语义和局部结构。低分时序主干利用较连续的年内观测，
高分逐景编码器提取各自原生尺度上的空间信息；日期与质量感知聚合适配不规则观测，
统一解码器在有支持的位置学习融合。上述分工构成本项目的主方案。

5m多光谱与2m PAN具有不同的光谱信息、计算成本和质量条件，因此分别设置编码与支持mask。
PAN先提取结构再按地面交叠聚合到5m，低分特征插值到同一网格，三路共同优化最终64维表征。
固定几何对齐使网格含义可核验；留出观测预测提供训练信号；空间独立的冻结嵌入评测决定方案是否有效。

\note{方法采用分为两层：主方案固定三分支接口、年度聚合和一次联合解码；多尺度骨干、学习上采样与冻结教师潜在目标
作为单项消融候选。候选方法只有在独立验证中证明收益、并满足8卡资源预算后才进入正式配置。
参考文献的作用是说明具体设计依据；组合后的效果仍需本项目实验验证。}

\subsection{数据、目标与结构的投入顺序}
AEF使用Landsat的互补光学与热红外观测，并结合观测重建、子集一致性和额外语义目标；
因此本地只有来源重建与球面瓶颈的实现，不能由STP结构相近推断语义能力相当\cite{aef}。
Galileo区分局部和全局学习，OlmoEarth也报告朴素潜在目标存在坍塌风险；这支持同时检查
语义、边界、有效秩与分组泛化，而非只优化重建值\cite{galileo,olmoearth}。
本地均匀性是场景均值上的Wang--Isola对数指数距离目标，AEF文中的batch正交目标采用
不同公式；即使权重都写0.05，也不能推断约束强度一致。调整该权重须另做消融，
本轮为保证可归因性保持既有值。
TESSERA最新v7（2026-04-12）仍只以S1/S2学习，强调全球打散与mixup，并区分约46M的
推理编码器和约1.34B参数的训练投影头\cite{tessera}。这说明Landsat不是学习有效年度语义的
必要输入，不能把本地差距全部归因于它的缺失。当前DDP已在每年的全国清单内随机打散，
不能照搬该论文的打散消融增益。其原始时序、10m/128维输出与本地月合成、5m/64维合同不同；
大投影头也占用训练资源，不能只比较推理参数。这些参考方向均未进入本轮冻结ABC配方。
局部结构还可考虑关系约束：DINOv3针对长训中密集特征退化提出Gram anchoring\cite{dinov3}。
若本地粗尺度语义约束持续伴随道路回退，可另测只在真实可见高分区域保留局部特征关系，
并与仅增加像素重建权重区分；其在本项目的效果尚无实验结论。
AEF的地形、LiDAR、环境量与部分地图/文本主要提供训练目标，而非全部作为推理输入；
因此不能把方法差距简化为本地少接了几路输入。若增加语义地图目标，必须另备未用于监督的
验证标签，不能继续把同一地图当独立评测。其精确vMF参数与本地维度缩放的高斯投影参数
也不是同一定义，应比较实际角噪声，不能直接按配置中的$\kappa$数值倍数修改。
两份实际checkpoint的本地$\exp(\log\kappa)$均为10，年度模型将其固定；
每维高斯标准差为$1/(10\sqrt{63})$，投影后小角近似对应精确vMF集中度约6{,}300。
固定种子的10万次CPU采样得到本地角偏移RMS约$5.71^{\circ}$，
而64维精确vMF、$\kappa=8{,}000$约$5.08^{\circ}$。
这说明两者实际扰动尺度接近，不能把配置中的10与8{,}000直接解释为800倍正则强度差；
该分布诊断没有改变模型，也不提供噪声最优值的实验结论。
测量条件也是目标设计的一部分：当前年度\texttt{DateDecoder}以日期条件化来源读出，
尚未使用太阳角、观测角或SAR轨道等逐观测条件。若这些变化被迫由嵌入承担，可能削弱
地表语义的稳定性；补充条件读出须先恢复可靠元数据，再做独立消融，不能由重建变好直接推断收益。

本地722{,}118条Landsat候选因云影QA缺失被明确排除，两版已训练模型均未使用它。
接入前需恢复原始产品身份和QA，区分TOA与地表反射率、光学与温度的单位及scale/offset，
处理饱和、填充、云影，并只用train重算统计量；USGS的QA\_PIXEL与QA\_RADSAT给出
相应质量标志\cite{landsatqa}。Landsat可补时间覆盖、光谱和热信息，不能把30m/100m观测
上采样后当作5m真值。2020年不存在Landsat 9观测，新增年份与卫星须保持真实时间合同。
当前物化与STAC接口只取六个光学波段和\texttt{QA\_PIXEL}，年度注册的Landsat为六通道，
没有读取热红外。恢复QA并不会自动增加温度信息；热红外须单独登记波段、单位、质量与目标头。
两版现有权重也未包含Landsat的stem与目标头。扩展source会改变参数及来源合同，
须使用新的训练配置和经审计的新增模块初始化；严格恢复不会把新增模态当作无变化的resume。
还须区分QA的用途：AEF文中相关质量mask用于训练，推理不向模型提供该mask；本地v6则
已在输入上保守屏蔽云影与无效像元。两条流水线的保留覆盖、输入分布和遮挡超参数需分别核验。
同样须先恢复已有S2的输入合同：归档十通道样本抽查没有波段描述，v6仍登记为
\texttt{requires\_provenance\_review}；当前STAC/物化接口选择的是十一光学波段加SCL。
两者不能未经登记直接互换，应依据原始产品证据核验顺序、单位与处理级别，
不能凭常见波段顺序或像元强度猜测。该抽查是元数据诊断，不证明所有归档文件均缺少元数据。

增加数据应先改善互补性：更多独立位置、生态区、年份和稀缺高分支持，比重复现有epoch
更有机会补足语义与细节缺口，但收益仍需对照验证。针对5m目标，还需要可靠的建筑/道路
标注与定位核验，避免OSM缺失被当作负样本。结构方面优先核验高分编码、分支尺度和
监督梯度，再考虑扩大高分骨干或多尺度解码；当前100M已经表明仅扩大低分STP不足以保证收益。
上述是结合文献与本地证据的投入判断，不是尚未实施的数据或结构改动的性能承诺。
AEF公开年度嵌入也可作为另一种10m教师目标候选；这需要先为训练位置建立可追溯缓存，
并另做外部教师消融。本轮C使用本地4M教师，没有把AEF产品接入监督，也不把其10m插值当5m真值。
采样比例还须说明地理基数：全国1\%与全球约1.1\%并不等价；月合成、原始景、父格和
像元的计数单位也不同。应记录独立位置、年份、有效观测与质量覆盖，避免用相同epoch或相近比例
掩盖实际训练信息量的差异\cite{aef}。

AEF论文v2的生产修订说明同样值得参考：其v2.1放宽必须具备特定来源或目标的采样条件，
修正年初时间编码与推理误用训练抽帧的问题，并调整地图监督混合和多分辨率重网格方式。
论文主结果对应AEF v2.0，不能与修订后生产产品混称。对本项目的启发是把来源缺测、
时间边界、训练与导出的输入策略、网格相位列为实测合同；这属于从文献推导的审计方向，
并不意味着本地存在相同缺陷，也不构成本轮额外训练组。

\subsection{下一轮投入的顺序与验收条件}
2026-10-16进一步确认已有独立目标数据包后，下一阶段优先采用“修复训练与输入合同，
接入DEM和WorldCover的训练期监督”路线，研究结论与第三版模型见第14节。
冻结4M蒸馏降为可选研究对照；下表不把教师预测头调参作为必经步骤。

\noindent\begin{tabularx}{\linewidth}{@{}p{2.8cm}YY@{}}
\toprule 顺序 & 具体投入 & 进入下一步的条件\\\midrule
先补可验证合同 & 核实S2波段与单位；独立核验高分定位；补多地理块的建筑、道路真值 & 能区分真实细节变化与OSM漏标、几何误差及头随机性\\
先接已有监督目标 & DEM高程、WorldCover分别与联合对照；CLCD先保留作跨产品诊断 & 目标只进入监督分支，完成有效覆盖、几何、尺度、年份与划分验收\\
再改目标与优化配方 & 固定观测重建与高分可见监督，登记辅助头初始化、学习率及损失梯度 & 同样本、同更新对照并报告算力差异，语义及细节分项不过度退化\\
再增加互补数据 & 新位置、生态区与年份；优先补缺少的有效5m多光谱，Landsat按QA及光学/热目标分阶段接入 & 按唯一位置、有效观测和高分支持计信息增量；保持未用于监督的空间留出\\
再给高分路径容量 & 固定现有低分STP，对高分编码或多尺度解码作单项扩容 & 实测吞吐、显存与有高分位置的边界收益，排除仅5m输出网格带来的错觉\\
最后扩大长训 & 只延长已通过对照的配方，记录新增数据与计算量 & 多个继续训练种子和空间确认支持结论；正式收益仍需独立预训练重复与真实5m验证\\\bottomrule
\end{tabularx}
本表是实验的优先级；独立目标的研究结论与v3结果见第14节。

\reportpage
\section{独立目标监督研究与第三、四版模型}
\label{sec:targets-next-plan}
\note{\status{本节范围}：2026-10-16至10-19在原100M上进行的独立目标监督研究（原v1.10--v1.27第14节）
已压缩为下列结论；逐轮研究记录、执行合同和中间制品保留在Git标签\texttt{archive/annual-research-20261019}
（提交\texttt{07da328}）中，不再在本正文维护。原100M权重与\texttt{v2.0}标签不变。}

\subsection{研究阶段的结论}
\textbf{问题。} 原100M只依靠观测重建与均匀性训练，没有静态地理监督。本阶段检验：在训练期加入
与输入分离、只进入损失的目标，能否改善冻结嵌入的语义，同时保留5m细节。
目标共七项：DEM高程与坡度、WorldCover、CLCD、VIIRS夜光、OSM 10m语义与2.5m结构
\cite{copdem,worldcover,clcd,vnl}。它们属于四类证据族，不是七套统计独立的真值；
推理仍只读取S1/S2及可用高分观测。

\textbf{期间发现并修复的实现缺陷。}
\begin{itemize}
\item \textbf{PAN可见掩膜：}旧覆盖判定误删48.6\%的有效PAN像元；改为精确地面覆盖（\path{annual_exact_mask_coverage}）。
\item \textbf{遮挡区上下文梯度：}被遮挡像元的上下文通路梯度为零；改为上下文池化（\path{annual_context_pooling}）。
\item \textbf{多PAN参数身份：}四个PAN投影曾按Python集合顺序注册，DDP广播与优化器槽位可能静默错配；
现按配置顺序注册，DDP包装前跨rank核对参数名称/形状，checkpoint保存有序参数名\cite{torchddp26,torchoptimizer26}。
\item \textbf{读出数值波动：}同一特征的卷积读出在CUDA上非逐位可复现；城市读出增加确定性模式并记录数值策略\cite{torchdet26}。
\end{itemize}
第11节涉及旧多PAN实现的多卡配方差值，因缺少逐rank身份证明，只保留为制品分数，不作因果归因。

\textbf{续训配方筛选。} 所有候选均从原100M续训1,024或3,116步，双种子配对，用同一原始参考和固定读出评估。
\begin{itemize}
\item 七目标监督确实被学到：新空间256父位置上WorldCover mIoU和40m DSM误差在各配方中均显著改善。
\item 但普通续训会侵蚀5m道路细节：不加保持项的九组配方中，四个监督模型道路F1下降0.86--1.79个百分点，
重建控制组也同样下降，说明主因是续训本身而非新增监督。
\item 单独冻结vMF投影、降低学习率或降低监督权重都不能同时保住道路与语义增益。
\item 加入相对原100M的冻结特征保持项后道路得以保留。三种保持形式中，
10m区块加权保持（粗尺度系数1、细节系数10）被带入全范围3,116步配对，两种子相对原100M的道路等保护项均通过。
\end{itemize}

\begin{center}\small
\begin{tabularx}{\linewidth}{@{}lYrr@{}}
\toprule 阶段 & 配方 & 平均$\Delta S$ & 平均$\Delta F$\\\midrule
无保持、0.20总权重 & 七目标 vs 匹配控制 & +0.9--+1.6 & 不稳定；道路超出容差\\
整向量保持 & 七目标 vs 匹配控制 & $-0.006$ & +0.352\\
区块残差保持 & 七目标 vs 匹配控制 & +1.094 & $-0.036$\\
区块加权保持 & 七目标 vs 匹配控制 & +0.167 & +0.044\\
区块加权保持，全范围3,116步 & 七目标 vs 原100M & +0.142 & +0.132\\\bottomrule
\end{tabularx}
\end{center}
单位为百分点，均为两个续训种子的平均配对差，2,000次整空间组bootstrap；
原研究按逐种子正增量和保护条件判定，上述配方均未通过该严格门槛。
2026-10-19改为以“相对原100M有改进且道路不退”为目标，据此选定最后一行配方作为第三版长训配方。
\source{[E12]原100M修复与选择记录；[E15]全范围目标缓存与训练池；归档标签内的完整研究证据。}

\subsection{第三版（v3）：区块加权保持的8遍续训}
\textbf{配方。} 从原100M初始化，载入全范围预热的七个监督头；主干学习率$1\times10^{-5}$、头学习率×30，
vMF投影冻结；七目标总权重0.2；以原100M（修复后实现）为冻结参考，区块加权特征保持系数10；
开启PAN精确掩膜、上下文池化与梯度检查点。训练池为允许监督范围49,868位置、双年99,736条记录。

\textbf{训练。} 8张A100，批次$4\times8$，共8遍、24,928次更新、797,696次年度曝光；
墙钟58,425\,s（16.2\,h），13.65样本/s，单卡峰值29.0\,GB，非有限步数为0。
每遍保存checkpoint；评测前预先指定第8遍为最终模型，不按事后最好epoch选择。

\textbf{评测。} 与研究阶段完全相同的协议：124个弱标签语义位置（20.48\,km空间组分折，线性头）；
16城市203个5m候选（城市分折，600步线性与$3\times3$卷积头，确定性读出，读出种子偏移0和10）；
另做5m特征池化到10m的读出。重新计算的原100M参考与历史值逐位一致（$S=64.944$，原生读出$F=38.605$）。

\begin{center}\small
\begin{tabularx}{\linewidth}{@{}lrrY@{}}
\toprule 指标（两读出种子平均） & 原100M & v3第8遍 & 配对差及95\%区间（pp）\\\midrule
语义$S$ & 64.944 & 64.908 & $-0.04$ $[-0.59,+0.43]$\\
\quad 地表覆盖mIoU & 43.777 & 44.042 & $+0.27$ $[-0.04,+0.63]$\\
\quad 10m水体F1 & 86.112 & 85.773 & $-0.34$ $[-1.40,+0.44]$\\
5m细节$F$ & 38.478 & 39.015 & $+0.54$ $[+0.02,+1.08]$\\
\quad 道路F1 & 23.696 & 23.920 & $+0.22$ $[-0.40,+0.80]$\\
\quad 道路边界F1 & 33.932 & 34.145 & $+0.21$ $[-0.56,+1.01]$\\
\quad 5m水体F1 & 71.782 & 72.685 & $+0.90$ $[-0.04,+2.02]$\\
\quad 建筑F1 & 24.500 & 25.310 & $+0.81$ $[-1.04,+2.28]$\\
池化到10m后的$F$（单读出种子） & 37.906 & 38.648 & $+0.74$ $[+0.10,+1.51]$\\\bottomrule
\end{tabularx}
\end{center}

\begin{figure}[h]\centering
\includegraphics[width=\linewidth]{/data/heyuhang/xuannv_embedding/outputs/v3/eval/figures/v3_epoch_deltas.pdf}
\caption{v3各遍相对原100M的配对差（百分点）。阴影为整空间组95\%区间，菱形为预先指定的第8遍终点；
$S$与地表覆盖为单一线性读出，$F$与道路F1为两个读出种子平均。}
\end{figure}

\textbf{解读。}
\begin{itemize}
\item \textbf{道路细节保住了。}这是此前所有续训配方的失败点；v3的道路与道路边界F1均不低于原100M。
\item \textbf{5m细节小幅提升}，两读出种子平均$+0.54$个百分点，区间下界略高于零；池化到10m后同样为正。
\item \textbf{语义持平。}地表覆盖增益随训练单调衰减（第1遍$+0.60$到第8遍$+0.27$），10m水体略降；
一个可能原因是保持项持续把特征拉回原模型，抵消了独立监督带来的语义变化，这一点尚未单独验证。
\item 第1遍的细节分数最高，但它是事后观察；按预先规定报告第8遍。
\item 延长训练没有带来随遍数稳定增长的收益；其原因（保持强度、学习率日程或目标信息量）未作单独隔离。
\end{itemize}
124个语义位置中110个已被原预训练访问；5m标签来自OSM弱标签，道路、建筑、水体实际计分支持为186、38、80个patch。
以上是相对原100M的稳中小幅改进，不是独立人工真值下的精度认证。
\source{[E65]v3训练配置、日志与权重；[E66]v3评测计划、逐遍读出、配对bootstrap脚本与结果。}

\subsection{第四版（v4）：修复后从头预训练}
\textbf{设置。} v3被保持项锚定在原100M上，无法回答两处修复对预训练本身的影响。v4使用与原100M完全相同的配置
（数据、200遍、学习率$1\times10^{-4}$、编译、批次32、种子42），只打开
\path{annual_context_pooling}与\path{annual_exact_mask_coverage}，不加独立目标或保持项。
2026-10-22至10-26在8张A100上完成632{,}000次更新，墙钟398{,}008\,s（4.6天），50.8样本/s，非有限步数为0。
预先指定第200遍为终点，另评第50/100/150遍以观察趋势。评测协议与v3完全相同。

\begin{center}\small
\begin{tabularx}{\linewidth}{@{}lrrrY@{}}
\toprule 指标（两读出种子平均） & 原100M & v3第8遍 & v4第200遍 & v4配对差及95\%区间（pp）\\\midrule
语义$S$ & 64.944 & 64.908 & 61.908 & $-3.04$ $[-8.71,+1.67]$\\
\quad 地表覆盖mIoU & 43.777 & 44.042 & 42.227 & $-1.55$ $[-4.26,+1.60]$\\
\quad 10m水体F1 & 86.112 & 85.773 & 81.590 & $-4.52$ $[-14.82,+3.32]$\\
5m细节$F$ & 38.478 & 39.015 & 36.560 & $-1.92$ $[-3.28,-0.62]$\\
\quad 道路F1 & 23.696 & 23.920 & 22.063 & $-1.63$ $[-3.90,+0.82]$\\
\quad 道路边界F1 & 33.932 & 34.145 & 32.476 & $-1.46$ $[-4.26,+1.10]$\\
\quad 5m水体F1 & 71.782 & 72.685 & 69.063 & $-2.72$ $[-7.78,+2.02]$\\
\quad 建筑F1 & 24.500 & 25.310 & 22.637 & $-1.86$ $[-4.50,+1.36]$\\
池化到10m后的$F$ & 37.906 & 38.648 & 36.120 & $-1.79$ $[-3.41,+0.52]$\\\bottomrule
\end{tabularx}
\end{center}

\begin{center}\small
\begin{tabular}{@{}lrrrr@{}}
\toprule v4遍数 & 50 & 100 & 150 & 200\\\midrule
语义$S$ & 61.61 & 62.00 & 61.25 & 61.91\\
$\Delta F$（两读出种子，pp） & $-0.24$ & $-0.45$ & $-1.37$ & $-1.92$\\\bottomrule
\end{tabular}
\end{center}

\textbf{解读。}
\begin{itemize}
\item \textbf{v4整体劣于原100M。}5m细节下降的区间完全低于零；语义平均下降3个百分点，区间跨零。
v4是独立训练的模型，与参考的配对相关性低，因此区间比v3宽得多。
\item 劣势在第50遍已出现，细节差值随训练继续扩大，说明问题来自配方而非训练不足。
\item \textbf{训练损失更低、表征更差。}训练末200次记录的平均重建损失为0.21，原100M为0.45；
同一批40个城市位置上，嵌入有效秩为25.4（原100M 30.9、v3 30.7），5m高频能量占比为1.9\%（原100M 2.3\%）。
一个与现象相容的解释是：上下文池化让遮挡像元的嵌入直接承担遮挡值预测，重建变容易，
嵌入不再需要编码丰富的地物信息。该机制未经单独验证；两项修复同时打开，也无法区分各自的作用。
\item v4与原100M各只有一个预训练种子，从头训练的种子间方差未知；但细节分数、有效秩与高频占比方向一致。
\end{itemize}
v4不作为候选版本。两项修复在v3中与独立目标、保持项同时使用，未观察到细节退化；但单独用于从头训练时整体变差，
因此不能把它们一律视为改进。
\source{[E67]v4配置、日志与逐遍权重；[E66]同一评测计划与\texttt{results\_v4.json}、\texttt{two\_seed\_results\_v4.json}。}

\subsection{下一步}
\begin{enumerate}
\item 以v3第8遍为当前候选；发布前需单独确认。
\item 若要定位v4的退化来源，用短预算或4M规模做单因素对照：只开精确PAN掩膜、只开上下文池化。
\item 若目标转向语义，应调整保持项的强度或形式，而不是延长训练。
\item 原S2生产数值链、Landsat云质量QA与独立人工5m真值缺口仍未关闭\cite{landsatqa}。
\end{enumerate}

\reportpage
\phantomsection
\addcontentsline{toc}{section}{参考文献}
\begin{thebibliography}{99}\small
\bibitem{aefgcs} Google. AlphaEarth Foundations GCS data. 2026-07-07更新，本版2026-10-18实际核对。\url{https://developers.google.com/earth-engine/guides/aef_on_gcs_readme}.
\bibitem{aefcatalog} Google. Satellite Embedding V1. 本版2026-10-18实际核对。\url{https://developers.google.com/earth-engine/datasets/catalog/GOOGLE_SATELLITE_EMBEDDING_V1_ANNUAL}.
\bibitem{aef} Christopher F. Brown et al. AlphaEarth Foundations: An embedding field model for accurate and efficient global mapping from sparse label data. 2025. \url{https://arxiv.org/abs/2507.22291}. 本轮升级核对v2的S15、S16、S21与S22.3；历史实现并非完整复现。
\bibitem{anysat} Guillaume Astruc et al. AnySat: One Earth Observation Model for Many Resolutions, Scales, and Modalities. CVPR, 2025. \url{https://arxiv.org/abs/2412.14123}.
\bibitem{tessera} Zhengpeng Feng et al. TESSERA: Temporal Embeddings of Surface Spectra for Earth Representation and Analysis. 2025--2026. \url{https://arxiv.org/abs/2506.20380}. 历史方案核对v1；本轮调研另核对2026-04-12的v7第3、5节与附录A.2、A.7。
\bibitem{barlow} Jure Zbontar et al. Barlow Twins: Self-Supervised Learning via Redundancy Reduction. ICML, 2021. \url{https://proceedings.mlr.press/v139/zbontar21a.html}.
\bibitem{fpn} Tsung-Yi Lin et al. Feature Pyramid Networks for Object Detection. CVPR, 2017. \url{https://arxiv.org/abs/1612.03144}.
\bibitem{swin} Ze Liu et al. Swin Transformer: Hierarchical Vision Transformer using Shifted Windows. ICCV, 2021. \url{https://arxiv.org/abs/2103.14030}.
\bibitem{scalemae} Scale-MAE: A Scale-Aware Masked Autoencoder for Multiscale Geospatial Representation Learning. ICCV, 2023. \url{https://arxiv.org/abs/2212.14532}.
\bibitem{vmf} von Mises--Fisher Loss: An Exploration of Embedding Geometries for Supervised Learning. 2021. \url{https://arxiv.org/abs/2103.15718}.
\bibitem{dysample} Wenze Liu et al. Learning to Upsample by Learning to Sample. ICCV, 2023. \url{https://arxiv.org/abs/2308.15085}.
\bibitem{featup} Stephanie Fu et al. FeatUp: A Model-Agnostic Framework for Features at Any Resolution. ICLR, 2024. \url{https://arxiv.org/abs/2403.10516}.
\bibitem{ijepa} Mahmoud Assran et al. Self-Supervised Learning from Images with a Joint-Embedding Predictive Architecture. CVPR, 2023. \url{https://arxiv.org/abs/2301.08243}.
\bibitem{galileo} Gabriel Tseng et al. Galileo: Learning Global \& Local Features of Many Remote Sensing Modalities. 2025. \url{https://arxiv.org/abs/2502.09356}. 核对v3。
\bibitem{olmoearth} Henry Herzog et al. OlmoEarth: Stable Latent Image Modeling for Multimodal Earth Observation. 2025. \url{https://arxiv.org/abs/2511.13655}.
\bibitem{dinov3} Oriane Sim\'eoni et al. DINOv3. 2025. \url{https://arxiv.org/abs/2508.10104}.
\bibitem{landsatqa} USGS. Landsat Collection 2 Quality Assessment Bands. \url{https://www.usgs.gov/landsat-missions/landsat-collection-2-quality-assessment-bands}. 核验于2026-10-15。
\bibitem{copdem} Copernicus Data Space Ecosystem. Copernicus DEM: Global and European Digital Elevation Model.
\url{https://dataspace.copernicus.eu/explore-data/data-collections/copernicus-contributing-missions/collections-description/COP-DEM}.
产品定义、DSM范围及EGM2008米制基准核验于2026-10-16；本地原生派生链以[E15]为准，不作独立测量精度认证。
\bibitem{worldcover} ESA WorldCover. Data access: WorldCover 2020 and 2021.
\url{https://esa-worldcover.org/en/data-access}. 核验于2026-10-16；本地类别、覆盖及有效mask以[E13]登记为准。
\bibitem{s2products} ESA / Copernicus. S2 Products, SentiWiki.
\url{https://sentiwiki.copernicus.eu/web/s2-products}. 核验于2026-10-16；物理表面反射率须结合产品元数据中的BOA偏移和量化因子。
\bibitem{s2processing} ESA / Copernicus. S2 Processing, historical archive reprocessing, SentiWiki.
\url{https://sentiwiki.copernicus.eu/web/s2-processing}. 核验于2026-10-16；历史成像年份不能替代实际处理基线证明。
\bibitem{clcd} Jie Yang and Xin Huang. The 30m annual land cover dataset and its dynamics in China from 1990 to 2019.
Earth System Science Data, 2021. \url{https://doi.org/10.5194/essd-13-3907-2021}.
核验于2026-10-16；作为CLCD方法与30m来源背景，2020/2021本地文件身份见归档标签中的研究证据。
\bibitem{vnl} Earth Observation Group. VIIRS Nighttime Light: Annual VNL.
\url{https://eogdata.mines.edu/products/vnl/}. 核验于2026-10-16；本地为已发布年度平均值派生标签，
源文件名、15角秒网格及转换范围见归档标签中的研究证据。
\bibitem{salt} Xianhang Li et al. Rethinking JEPA: Compute-Efficient Video SSL with Frozen Teachers. 2025. \url{https://arxiv.org/abs/2509.24317}.
\bibitem{presto} Gabriel Tseng et al. Lightweight, Pre-trained Transformers for Remote Sensing Timeseries. 2023. \url{https://arxiv.org/abs/2304.14065}.
\bibitem{ocm} Nick Wright. OmniCloudMask, 作者实现及模型。\url{https://github.com/DPIRD-DMA/OmniCloudMask}. 本地运行版本及权重摘要以[E2]为准，不以仓库最新状态反推已运行版本。
\bibitem{gfguide} SpaceWill. GaoFen-1 Satellite Imagery User Guide. 产品说明；用于传感器与波段惯例，不证明当前切片的处理精度。 \url{https://en.spacewillinfo.com/uploads/soft/210106/8-210106153534.pdf}.
\bibitem{alphaearth} Christopher F. Brown et al. AlphaEarth Foundations: An embedding field model for accurate and efficient global mapping from sparse label data. 2025. \url{https://arxiv.org/abs/2507.22291}. 本报告仅借鉴年度嵌入和条件时空汇总方向，不使用其嵌入作为本地监督标签。

\bibitem{lpft} Ananya Kumar, Aditi Raghunathan, Robbie Jones, Tengyu Ma, Percy Liang.
Fine-Tuning can Distort Pretrained Features and Underperform Out-of-Distribution. ICLR, 2022.
\url{https://arxiv.org/abs/2202.10054}.
\bibitem{pcgrad} Tianhe Yu et al. Gradient Surgery for Multi-Task Learning. NeurIPS, 2020.
\url{https://proceedings.neurips.cc/paper_files/paper/2020/hash/3fe78a8acf5fda99de95303940a2420c-Abstract.html}.
\bibitem{gradnorm} Zhao Chen et al. GradNorm: Gradient Normalization for Adaptive Loss Balancing in Deep Multitask Networks. ICML, 2018.
\url{https://proceedings.mlr.press/v80/chen18a.html}.
\bibitem{torchadamw26} PyTorch contributors. AdamW implementation, v2.6.0.
\url{https://raw.githubusercontent.com/pytorch/pytorch/v2.6.0/torch/optim/adamw.py}.
\bibitem{torchddp26} PyTorch contributors. DistributedDataParallel implementation and parameter-order contract, v2.6.0.
\url{https://raw.githubusercontent.com/pytorch/pytorch/v2.6.0/torch/nn/parallel/distributed.py}.
\bibitem{torchoptimizer26} PyTorch contributors. Optimizer state dictionary and positional restoration, v2.6.0.
\url{https://raw.githubusercontent.com/pytorch/pytorch/v2.6.0/torch/optim/optimizer.py}.

\bibitem{torchrepro26} PyTorch contributors. Reproducibility, v2.6.0.
\url{https://raw.githubusercontent.com/pytorch/pytorch/v2.6.0/docs/source/notes/randomness.rst}.
\bibitem{torchnumerics26} PyTorch contributors. Numerical accuracy, v2.6.0.
\url{https://raw.githubusercontent.com/pytorch/pytorch/v2.6.0/docs/source/notes/numerical_accuracy.rst}.

\bibitem{lwfretention} Li Z, Hoiem D. Learning without Forgetting. arXiv:1606.09282v3, 2017.
\url{https://arxiv.org/abs/1606.09282v3}.
\bibitem{lflretention} Less-forgetting Learning in Deep Neural Networks. arXiv:1607.00122v1, 2016.
\url{https://arxiv.org/abs/1607.00122v1}.

\bibitem{podnet2020} Douillard A, Cord M, Ollion C, Robert T, Valle E. PODNet: Pooled Outputs Distillation for Small-Tasks Incremental Learning. ECCV, 2020.
\url{https://www.ecva.net/papers/eccv_2020/papers_ECCV/papers/123650086.pdf}.
\bibitem{plop2021} Douillard A, Chen Y, Dapogny A, Cord M. PLOP: Learning Without Forgetting for Continual Semantic Segmentation. CVPR, 2021.
\url{https://openaccess.thecvf.com/content/CVPR2021/html/Douillard_PLOP_Learning_Without_Forgetting_for_Continual_Semantic_Segmentation_CVPR_2021_paper.html}.

\bibitem{torchdet26} PyTorch contributors. torch.use\_deterministic\_algorithms, PyTorch 2.6. \url{https://docs.pytorch.org/docs/2.6/generated/torch.use_deterministic_algorithms.html}.
\end{thebibliography}

\reportpage
\appendix
\section{本地证据登记与版本维护}
\subsection{证据优先级与快照}
本地状态使用实际产物和代码，方法论使用原论文，拟实施配置标明为设计。旧Markdown叙述与机器摘要冲突时，
先核验版本与时间，不以“最新”字样自动覆盖证据。本版完成PAN筛选、吉林一号质量筛选、全量高分相对配准、
跨集合重复审计、600个真实样本读取、年度模型单元测试、8卡4步真实数据短跑，以及200轮正式训练的收敛、
冻结嵌入探针评测、与外部年度嵌入基线在同一协议下的对照和5m原生协议评测；没有独立真值下的精度认证，
也没有生产推理条件的闭环结论。
以下路径以服务器 \texttt{/data/heyuhang/processed/} 为基准；完整SHA-256及源码文件列表保存在
\texttt{docs/report-evidence.json}，它是证据索引，不是第二份方案正文。

\begingroup\small\renewcommand{\arraystretch}{1.1}
\noindent\begin{tabularx}{\linewidth}{@{}p{.8cm}p{5.4cm}Y@{}}
\toprule 编号 & 文件或证据 & 本报告使用范围\\\midrule
E1 & \path{national_annual5m_quality_v6/summary.json} & 三路年度规模、支持计数、放行标识；不代表全量科学审计。\\
E2 & \path{national_annual5m_quality_v3/run.json} & S2处理完成状态、总数、模型版本及处理规则。\\
E3 & \path{national_annual5m_quality_v6/loader-check.json} & 六份清单各100次读取、逐来源原生尺寸；不推及未检查样本。\\
E4 & \path{national_training_v1/summary.json} & 原catalog与父格背景；状态字段不替代后续年度放行。\\
E5 & 历史基线模型、年度读取器、质检代码 & 保留原源码身份；当前新增监督实现另见\texttt{current\_source\_sha256}。\\
E6 & \path{national_annual5m_quality_v6/pan-screening-summary.json} & 同景MUX关联、云影传播、面积阈值和PAN筛选计数。\\
E7 & \path{national_annual5m_quality_v6/run.json} 及合并核验记录 & 5m波段映射、模型与处理器摘要、合并输入版本及最终合同核验。\\
E8 & \path{national_annual5m_quality_v6/registration-full-summary.json}、\path{pan-screening-summary.json}、\path{visual-review.json} & 全量相对配准、PAN筛选和目视诊断；用于说明隔离边界，不替代独立精度验证。\\
E9 & \path{outputs/annual5m/run_200ep_20260921_092149/}（\texttt{train.log}、\texttt{annual5m\_v1.pt}、\texttt{git\_sha.txt}） & 200轮正式训练日志、权重摘要与来源SHA；不含独立精度结论。\\
E10 & \path{outputs/eval/quick\_rapid124/}、\path{outputs/eval/city5m/} 结果与脚本 & 历史4M冻结探针、外部基线对照与逐层测量；弱标签，非独立真值。\\
E11 & \path{outputs/annual5m/run_100m_20261009_024811/} 配置、日志与权重 & 原100M参数、200轮训练、632{,}000次更新及来源核验。\\
E12 & \path{outputs/annual_abc/20261015_exact/}；诊断原件位于\path{20261015_screen/} & 梯度/PAN修复、十一组候选与三组种子复验、30份主/参考头结果、空间确认及真实8卡EMA恢复；共同重建目标不混用旧掩膜损失。\\
E13 & \path{outputs/target_inventory/20261016/summary.json}；\path{/data/wwj/xuannv_china_ai_ready_dataset/targets/} & 初次独立目标资产、位置匹配与静态数组抽样；保留原盘点范围。\\
E15 & \path{outputs/annual_targets/ready_20261016/}、\path{prepare_20261016/} & 原生DEM分卷SHA/CRC、40m缓存、冻结评测及嵌套训练清单、scope训练统计与S2追溯；未认证独立5m精度。\\
E65 & \path{outputs/v3/run_20261019_041538/}：配置、日志、逐遍权重与\path{git_sha.txt} & v3配方、24,928次更新、吞吐与显存；源码\texttt{3310907}。\\
E66 & \path{outputs/v3/eval/}：评测计划、分析脚本与结果JSON & v3逐遍读出及相对原100M的配对区间；参考读出复用原研究固定结果。\\
E67 & \path{outputs/v4/run_20261022_045453/}：配置、日志与逐遍权重 & v4训练记录；评测结果见E66中的\path{results_v4.json}。\\\bottomrule
\end{tabularx}
\endgroup
E14、E16--E64为研究阶段证据，路径已随服务器清理删除，完整条目保留在归档标签\path{archive/annual-research-20261019}的\path{docs/report-evidence.json}中。

\subsection{唯一正文的维护方法}
\begin{enumerate}
\item 修改本 \texttt{xuannv\_annual5m\_report.tex}，同时更新版本、核验日期及受影响的证据/参考文献。
\item 数量或状态变化时重新读取对应版本摘要，更新证据索引和图表；不得只改报告中的结果数字。
\item 执行 \texttt{make -C docs report}，检查日志、引用、中文复制与图表边界。仓库不维护并行方案正文；历史材料仅保留外部归档。
\end{enumerate}

PDF是阅读与分发格式；可持续编辑的主文件是TeX。生成PDF不作为Git源码提交，
以遵守仓库二进制制品策略；实际交付文件仍保留在本地。Tectonic/XeLaTeX依赖和缓存不进入仓库。

\noindent\begin{tabularx}{\linewidth}{@{}p{1.6cm}p{2.4cm}Y@{}}
\toprule 版本 & 日期 & 变更\\\midrule
v1.0--v1.4 & 2026-09-17/18 & 建立年度5m三分支方案；接入2m PAN；完成v6数据质量筛选、配准诊断与放行。\\
v1.7--v1.8 & 2026-09-20/10-02 & 年度模型实现与8卡短跑；4M 200轮训练、冻结嵌入评测与外部基线对照。\\
v1.9 & 2026-10-15 & 100M审计：两项实现缺陷修复、十一组短程配方对照与一次性选定。\\
v1.10--v1.27 & 2026-10-16/19 & 独立目标监督研究（七目标、多种保持形式、全范围3,116步配对）；逐轮记录见归档标签。\\
v3.0 & 2026-10-22 & 精简第14节；加入v3训练与评测、v4启动；证据索引只登记仍保留的制品。\\
v3.1 & 2026-10-27 & 加入v4结果（不作为候选）；第11节压缩为结论。\\\bottomrule
\end{tabularx}
\end{document}
